\documentclass{article}

\PassOptionsToPackage{numbers, compress}{natbib}
\usepackage[eandd, final]{neurips_2026}

\makeatletter\renewcommand{\@trackname}{40th Conference on Neural Information Processing Systems (NeurIPS 2026), The Predictive AI Evaluation Competition at NeurIPS 2026.}\makeatother
\usepackage[utf8]{inputenc}
\usepackage[T1]{fontenc}
\usepackage{hyperref}
\usepackage{url}
\usepackage{booktabs}
\usepackage{array}
\usepackage{amsfonts}
\usepackage{amsmath}
\usepackage{microtype}
\usepackage{xcolor}
\usepackage{graphicx}
\usepackage{multirow}
\usepackage{float}
\setlength{\textfloatsep}{12pt plus 2pt minus 4pt}\setlength{\floatsep}{10pt plus 2pt minus 2pt}
\usepackage{tcolorbox}\tcbuselibrary{breakable,skins}
\newtcolorbox{chg}[1]{breakable,enhanced,colback=red!14,colframe=red!14,boxrule=0pt,arc=0pt,left=2pt,right=2pt,top=1.5pt,bottom=1.5pt,before skip=2pt,after skip=2pt,before upper={\textcolor{red!70!black}{\small\bfseries[#1]}\ }}
\newtcolorbox{oldtxt}[1]{breakable,enhanced,colback=red!14,colframe=red!14,boxrule=0pt,arc=0pt,left=2pt,right=2pt,top=1.5pt,bottom=1.5pt,before skip=2pt,after skip=2pt,before upper={\textcolor{red!70!black}{\small\bfseries[S#1 was]}\ }}
\newtcolorbox{newtxt}[1]{breakable,enhanced,colback=green!14,colframe=green!14,boxrule=0pt,arc=0pt,left=2pt,right=2pt,top=1.5pt,bottom=1.5pt,before skip=2pt,after skip=2pt,before upper={\textcolor{green!45!black}{\small\bfseries[S#1 now]}\ }}
\newtcolorbox{addnote}[1]{breakable,enhanced,colback=green!14,colframe=green!55!black,boxrule=0.6pt,arc=1pt,left=3pt,right=3pt,top=2pt,bottom=2pt,before skip=4pt,after skip=4pt,fontupper=\small\itshape,before upper={\textcolor{green!45!black}{\upshape\bfseries NEW [#1]:}\ }}
\newtcolorbox{qnote}[1]{breakable,enhanced,colback=green!14,colframe=green!55!black,boxrule=0.6pt,arc=1pt,left=3pt,right=3pt,top=2pt,bottom=2pt,before skip=4pt,after skip=4pt,fontupper=\small,before upper={\textcolor{green!45!black}{\bfseries QUESTION [#1]:}\ }}
\newtcolorbox{fixnote}[1]{colback=red!14,colframe=red!55!black,boxrule=0.6pt,arc=1pt,left=3pt,right=3pt,top=2pt,bottom=2pt,before skip=4pt,after skip=4pt,fontupper=\small,before upper={\textcolor{red!70!black}{\bfseries TO FIX (needs Person-B's plot script) [#1]:}\ }}
\newtcolorbox{improve}[1]{colback=red!14,colframe=red!55!black,boxrule=0.6pt,arc=1pt,left=3pt,right=3pt,top=2pt,bottom=2pt,before skip=4pt,after skip=4pt,fontupper=\small,before upper={\textcolor{red!70!black}{\bfseries STILL TO IMPROVE [#1]:}\ }}
\newcommand{\oldin}[1]{\colorbox{red!14}{#1}}
\newcommand{\newin}[1]{\colorbox{green!14}{#1}}
\setlength{\textfloatsep}{10pt plus 2pt minus 2pt}\setlength{\floatsep}{8pt plus 2pt minus 2pt}\setlength{\intextsep}{8pt plus 2pt minus 2pt}\setlength{\abovecaptionskip}{4pt}

\newcommand{\minus}{\ensuremath{-}}
\newcommand{\D}[2]{\ensuremath{\Delta_{\mathrm{#1}-\mathrm{#2}}}}
\newcommand{\code}[1]{\texttt{#1}}
\newcommand{\ghlink}{\href{https://github.com/Agrover112/aims-neurips/tree/submission}{\textcolor[rgb]{0.05,0.25,0.65}{\underline{GitHub}}}}
\newcommand{\figdir}{figures/eandd}

\title{What Transfers to a New Benchmark? Priors, Not Item Estimates, for Label-Efficient LLM Evaluation}

\author{%
  Arshia Teymourtash\thanks{Work done equally.} \\
  KTH Royal Institute of Technology \\
  Stockholm, Sweden \\
  \texttt{arturt@kth.se}
  \And
  Ankit Grover$^{*}$ \\
  KTH Royal Institute of Technology \\
  Stockholm, Sweden \\
  \texttt{agrover@kth.se}
}

\begin{document}

\maketitle

\begin{abstract}
Evaluating an LLM on every item of every benchmark is expensive, so predicting item outcomes from a few revealed answers is attractive, but only if what is learned from past benchmarks carries over to new ones. We ask what does, in a controlled comparison where the models, the evaluation tasks and the order in which answers are revealed are held fixed and only the source of the prior changes. A prior learned from past benchmarks \emph{transfers} to a new benchmark. Estimates of individual item difficulty \emph{do not}. The prior's benefit is small on average and depends on the benchmark: it helps on benchmarks of \emph{typical} difficulty and hurts on much easier or much harder ones. A neural network that predicts item difficulty from item text adds little beyond the prior. Whether pooling answers across LLMs helps \emph{remains open}, because our evaluation reveals answers from one LLM per benchmark. Our entry to the Predictive AI Evaluation Competition (Codabench submission 1005155) is a Beta-Binomial predictor. We release the simulation code, the fixed evaluation samples and every saved score on \ghlink.
\end{abstract}

%==============================================================================
\section{Introduction}
%==============================================================================

Evaluating language models on every item of every benchmark is expensive, so a predictor that infers a model's per-item results from a handful of labelled answers is attractive. The natural recipe is to learn from past benchmarks how items behave and carry that to a new benchmark. Whether it carries is an empirical question, and the answer matters because new benchmarks and new models keep arriving. Most benchmark-prediction methods lose their advantage over random sampling when extrapolating to higher-performing models~\citep{zhang2026benchmark}, and amortized item-response models learn item difficulty from item text~\citep{truong2025amortized}. We ask what, among the things one can fit offline, actually transfers to an unseen benchmark.

We use the Predictive AI Evaluation Competition's setting: a \emph{task} is one subject (a model under given evaluation settings) answering one benchmark, labels are bought one at a time, and predictions are scored by Brier ALC at 0, 1, 3, 7, 15 and 31 labels. Our contributions:
\begin{itemize}
  \item \textbf{A controlled comparison.} Models, tasks and label order are fixed and only the source of the prior changes, in an online simulation of the competition protocol (Section~\ref{sec:models}).
  \item \textbf{What transfers.} A learned prior transfers to a new benchmark and item-level estimates do not. The prior helps on benchmarks of typical difficulty and hurts on much easier or harder ones, and item text adds almost nothing beyond it (Section~\ref{sec:results}).
\end{itemize}
The simulation code, the fixed evaluation samples and every saved score are on \ghlink.

%==============================================================================
\section{Related work}\label{sec:related}
%==============================================================================
\paragraph{Label-efficient evaluation.} Label-efficient methods estimate LLM performance from a small number of responses. tinyBenchmarks uses curated item subsets and item-response models to reproduce aggregate benchmark results~\citep{polo2024tinybenchmarks}. Work on predicting results for new models finds that most benchmark-prediction methods lose their advantage over random sampling when extrapolating to higher-performing models, while one estimator that combines predictions with random labels keeps it~\citep{zhang2026benchmark}. The shift there is in which models are evaluated, while ours is in which benchmark is evaluated. That work also estimates aggregate benchmark performance and not per-item outcomes, so we treat it as related evidence for our question and not as an answer to it. We predict on entirely held-out benchmarks, and our target is the probability that each unanswered item is answered correctly, given a few revealed labels. Active evaluation acquisition learns a policy that chooses which examples to label, from a candidate pool that stays visible throughout, and predicts the remaining outcomes from the chosen ones~\citep{li2025active}. Our labels arrive in a random order, so we do not study which items to label next. The question is what carries over to a held-out benchmark.

\paragraph{Item response theory.} It separates subject ability from item difficulty, with Rasch as its basic model~\citep{rasch1960}. Latent-space and multilingual extensions address interactions and language effects~\citep{jeon2021,lior2026}. Most directly related, a model trained on pooled responses predicts item difficulty from item text, so that new items need less calibration~\citep{truong2025amortized}. Concurrent work predicts the difficulty of agent tasks before the agent runs, from features such as text embeddings and the entropy of a generated reasoning trace. On unseen benchmarks the embeddings carry almost no signal and transfer is weak overall~\citep{krsteski2026difficulty}. It predicts a population-level difficulty and not a particular subject's outcome on each item. Text-predicted difficulty can assign parameters to previously unobserved questions, but its usefulness on an entirely new benchmark requires empirical validation. We compare simple probes of this kind (ridge regression and a small neural network on item embeddings, not the model above) with priors that carry over only the typical accuracy of a task, which separates information shared across benchmarks from item-level estimates. Item response theory gains most when many LLMs answer the same items, and our draws reveal answers from one LLM per benchmark, so sharing labels across LLMs is not tested.

\paragraph{Empirical Bayes and shrinkage.} Both carry population-level information to related tasks by learning a prior across them~\citep{robbins1956,efron1973,rao2015}. In our comparison, only the prior fitted on training benchmarks is empirical Bayes. The fixed Jeffreys prior~\citep{jeffreys1946} and the settings tuned on validation are separate prior sources.

\paragraph{Robust priors and prediction-powered inference.} Robust meta-analytic-predictive priors give a statistical precedent for mixture-based protection against prior--data conflict~\citep{schmidli2014robust}, although their clinical-trial setting and construction differ from ours. Prediction-powered inference combines predictions with observed labels, but targets valid inference for population quantities~\citep{angelopoulos2023}.

%==============================================================================
\section{Data and evaluation protocol}\label{sec:setting}
%==============================================================================
\label{sec:setting-online}
\paragraph{Data.} The data are a binary response matrix $Y\in\{0,1\}^{N\times M}$ with $N$ subjects (rows, index $i$) and $M$ items (columns, index $j$): $Y_{ij}=1$ if subject $i$ answers item $j$ correctly. Most entries are missing, since a subject answers only the items of the benchmarks it was evaluated on. Each item $j$ belongs to the dataset $d(j)$ of one of $D=17$ benchmarks, so a benchmark is a block of columns, and a \emph{task} is one row restricted to one block: subject $i$ on benchmark $d$.

Each item $j$ has an embedding $x_j\in\mathbb R^{1024}$ computed once from its text by the frozen BGE-M3 encoder~\citep{chen2024bge} with a context of 8,192 tokens, and normalised to unit length. We use BGE-M3 because the corpus is multilingual (it includes Chinese, Russian and several African languages) and some items are long, which English-only and 512-token encoders handle poorly.

\paragraph{Online simulation and splits.} The competition protocol is \emph{online}: for each task, labels are revealed one at a time and the predictor updates from them. Whatever the predictor knows before the first label is \emph{offline}, that is, fitted on the training benchmarks beforehand. Our \emph{online simulation} runs this protocol on recorded responses, with disjoint column blocks for training, validation, and test (11, 3, and 3 benchmarks): priors are fitted offline on training, and validation selects settings. The validation and test benchmarks are \emph{out of distribution} (OOD): their responses are excluded from offline fitting.

A \emph{competition-shaped draw} samples one subject per benchmark, so a predictor cannot borrow labels from other subjects on the same benchmark; this simulates the competition protocol on public data. As a separate \emph{in-distribution} (ID) evaluation, we hold out rows (subjects) within the training blocks, so its benchmarks and items have been seen during fitting, and retain the questions covered by the fit subjects. The offline comparator reuses its fitted benchmark rates and item difficulties where available, and its online update changes only the subject ability.

\paragraph{Budgets and scoring.} For each subject--benchmark pair we freeze a split of distinct items into acquisition and evaluation pools, assigning the extra item to acquisition when the count is odd. A frozen random acquisition order supplies nested prefixes of 0, 1, 3, 7, 15, and 31 labels. At each budget the predictor conditions on the allowed prefix and predicts the same evaluation items.

Evaluation outcomes are used only to score those predictions. We average Brier scores over query items, then average repeated appearances within each task and give tasks equal weight. Brier ALC combines the six budgets with weights $(0.1,0.2,0.2,0.2,0.2,0.1)$. Paired intervals resample whole tasks with their matched predictions.



%==============================================================================
\section{Methodology}\label{sec:models}
%==============================================================================
We first define the three predictors, then two ways of improving on them, adaptive priors and item-text probes, and finally the sources of the prior parameters that every predictor needs.

\subsection{Models}
The three predictors differ in what they model: a task's accuracy (Beta-Binomial), a subject's ability and each item's difficulty (Rasch), and the similarity between items (kernel Rasch).


\paragraph{Beta-Binomial.} Let $p$ be a task's accuracy, with a Beta prior $p\sim\mathrm{Beta}(\alpha,\beta)$. After $n$ revealed labels, $k$ of them correct, the posterior is $\mathrm{Beta}(\alpha+k,\beta+n-k)$ and the prediction, the same for every item of the task, is its mean
\begin{equation}\label{eq:bb}
\hat p=\frac{\alpha+k}{\alpha+\beta+n}.
\end{equation}
The prior acts like $\alpha+\beta$ extra labels, $\alpha$ of them correct. The Jeffreys prior is $\alpha=\beta=\tfrac12$~\citep{jeffreys1946}. This is the exact Bayesian sequential update; nothing is learned.

\paragraph{Rasch.} The probability that subject $i$ answers item $j$ correctly is
\begin{equation}\label{eq:rasch}
p_{ij}=\sigma(\theta_i-b_j),\qquad b_j=\delta_j-\mu_{d(j)},
\end{equation}
where $\sigma$ is the logistic function, $\theta_i$ is the subject's ability, and $b_j$ is the item's difficulty. This is the Rasch model~\citep{rasch1960} with the item difficulty split into a benchmark part and an item part: $\delta_j$ is the item's deviation from the average difficulty $-\mu_d$ of its benchmark $d=d(j)$, so a larger benchmark rate $\mu_d$ means an easier benchmark. The priors are $\mu_d\sim\mathcal N(m_0,\sigma_\mu^2)$, $\theta_i\sim\mathcal N(0,\sigma_\theta^2)$ and $\delta_j\sim\mathcal N(0,\sigma_\delta^2)$. At every budget we refit the maximum a posteriori (MAP) estimate from scratch on all labels revealed so far in the draw; an unlabelled item gets $\delta_j=0$. 
\paragraph{Kernel Rasch.} Kernel Rasch keeps Equation~\eqref{eq:rasch} and changes only the prior on the item deviations $\boldsymbol\delta=(\delta_1,\dots,\delta_M)^{\top}$. Rasch treats them as independent Gaussians, $\delta_j\sim\mathcal N(0,\sigma_\delta^2)$; kernel Rasch makes them jointly Gaussian, with covariance given by the similarity of the item embeddings:
\begin{equation}\label{eq:kernel}
\boldsymbol\delta\sim\mathcal N\big(\mathbf 0,\ K\big),\qquad K=\sigma_\delta^2\big[(1-\rho)\,I+\rho\,S\big],\quad S_{jj'}=\exp\!\Big(\frac{\cos(x_j,x_{j'})-1}{T}\Big),
\end{equation}
where $x_j$ are the item embeddings of Section~\ref{sec:setting}, $\rho\in[0,1]$ weights the similarity part $S$ and $T$ is its temperature. At $\rho=0$ the two priors coincide; for $\rho>0$, similar items share difficulty, while items of different benchmarks stay independent. After the refit, a labelled item keeps its fitted deviation, and an unlabelled item $j$ receives the posterior mean given the labelled items $L$ of its benchmark, $\hat\delta_j=\rho\,\sigma_\delta^2\,S_{jL}K_{LL}^{-1}\hat{\boldsymbol\delta}_L$, which is $0$ at $\rho=0$ as in Rasch. With one subject per benchmark, as in the competition-shaped draws, no labels are shared across subjects: the joint fit equals independent per-task fits (maximum difference $10^{-14}$ in our check), so Rasch there differs from the Beta-Binomial only in its prior.

\paragraph{Offline model that keeps its fitted values.} For the ID evaluation, the offline Rasch keeps its fitted $\mu_d$ and $\delta_j$ for a benchmark and items it has seen and updates only the subject ability $\theta_i$ online (a one-dimensional MAP on the task's own labels). For an unseen benchmark or item it falls back to $m_0$ and $0$, i.e.\ to the fitted prior alone.

\subsection{Adaptive priors}\label{sec:adaptive}
The predictors above keep their prior fixed whatever labels arrive. An \emph{adaptive prior} lets the first labels change it. We test three closed-form versions, each refit without any held-out benchmark.

\paragraph{Beta mixture.} The prior is a mixture of $K\in\{2,3\}$ Beta distributions with weights $w_c$, fitted by expectation maximisation on the accuracies of the training tasks. After $k$ of $n$ labels are correct, the weight of each component is updated by how well it explains them, and the prediction is the weighted mean of the components' posterior means,
\begin{equation}\label{eq:mixture}
w_c'\propto w_c\,\frac{\mathrm{B}(\alpha_c+k,\ \beta_c+n-k)}{\mathrm{B}(\alpha_c,\ \beta_c)},\qquad \hat p=\sum_c w_c'\,\frac{\alpha_c+k}{\alpha_c+\beta_c+n},
\end{equation}
where $\mathrm{B}$ is the Beta function.

\paragraph{Averaging over the prior.} Instead of one $(\alpha,\beta)$, we average the Beta-Binomial prediction over a grid of prior parameters, weighting each by the same factor $\mathrm{B}(\alpha+k,\beta+n-k)/\mathrm{B}(\alpha,\beta)$. We average over the prior's strength $\alpha+\beta$ at the fitted mean, or over its mean and strength together, in which case nothing is learned from training.

\paragraph{Switch to Jeffreys.} Once at least $k_0\in\{3,7\}$ labels are revealed, we predict with the Jeffreys prior if the accuracy so far, $k/n$, lies more than $\eta$ from the fitted prior's mean $\alpha/(\alpha+\beta)$, and with the fitted prior otherwise. We choose $\eta$ on validation.

\subsection{Probes}\label{sec:probes}
\label{sec:probes-text}
A probe predicts the difficulty of an item from its features alone, before any label of its benchmark is seen. The target is the mean binary correctness per item, centered within benchmark for fitting. We compare a ridge probe of item embeddings, a probe using log character count, and a sampled random score. We also train two small neural probes (MLPs) on the item embeddings, one on item pass rate and one on the fitted Rasch item difficulty, using 8,462 items from the 11 training benchmarks only.

We score a probe on each benchmark $d$ by two numbers, and average each over benchmarks. Let $y_j$ be the observed difficulty of item $j$ and $\hat y_j$ the probe's prediction. The \emph{Spearman correlation} $\rho_d=\operatorname{corr}\big(\operatorname{rank}(\hat y_j),\operatorname{rank}(y_j)\big)$ compares the two orderings of the items: it is $1$ if the probe orders them perfectly, $0$ if its order is unrelated to the truth, and $-1$ if it is reversed. The \emph{pairwise accuracy} $A_d=\frac{1}{|P_d|}\sum_{(j,k)\in P_d}s_{jk}$ is the share of item pairs the probe orders correctly, where $P_d$ is the set of pairs with $y_j\neq y_k$ (observed ties are excluded), $s_{jk}=1$ if $\hat y_j-\hat y_k$ has the same sign as $y_j-y_k$, $s_{jk}=\tfrac12$ if the probe predicts a tie, and $s_{jk}=0$ otherwise. A random score has $A_d=0.5$ on average. The ID diagnostic holds out items within the 11 training benchmarks; the OOD columns evaluate entire unseen benchmarks. All methods use matched items and sampled pairs.

To test whether a probe helps prediction, we give the item deviations of Rasch a prior mean taken from it,
\begin{equation}\label{eq:probe-prior}
\delta_j\sim\mathcal N\big(\gamma\,z_j,\ \sigma_\delta^2\big),
\end{equation}
where $z_j$ is the probe's predicted difficulty of item $j$, centred within its benchmark (the benchmark rate $\mu_d$ absorbs the average), and $\gamma$ is a scale; $\gamma=0$ is the fitted Rasch. The probes use item features only, never outcomes.



\subsection{Priors}
Every predictor above needs prior parameters, and where they come from is the one thing our comparison changes. We set them from one of three sources. \emph{Fixed in advance} means nothing is learned, so the Beta-Binomial uses Jeffreys and the Rasch models use unit widths. \emph{Fitted on training} means estimated from the 11 training benchmarks only. \emph{Tuned on validation} means the best setting on the three validation benchmarks from a grid. Table~\ref{tab:settings} lists the settings.

\begin{table}[H]
\centering
\caption{Prior parameters by model and source. The prior mean of the Beta-Binomial is $\alpha/(\alpha+\beta)$ and its strength is $\alpha+\beta$. Kernel Rasch keeps $\rho$ and $T$ from validation, so its kernel is not fitted on training.}
\label{tab:settings}
\footnotesize\setlength{\tabcolsep}{4pt}\renewcommand{\arraystretch}{1.15}
\begin{tabular}{@{}l>{\raggedright\arraybackslash}p{3.0cm}>{\raggedright\arraybackslash}p{4.1cm}>{\raggedright\arraybackslash}p{4.1cm}@{}}
\toprule
Model & Fixed in advance & Fitted on training & Tuned on validation\\
\midrule
Beta-Binomial & $\alpha=\beta=\tfrac12$ & $\alpha=1.92$, $\beta=1.76$ & $\alpha=0.4$, $\beta=0.6$\\
Rasch & $m_0=0$, $\sigma_\mu=\sigma_\theta=\sigma_\delta=1$ & $m_0=0.175$, $\sigma_\mu=0.75$, $\sigma_\theta=1.38$, $\sigma_\delta=1.49$ & $m_0=0.048$, $\sigma_\mu=2$, $\sigma_\theta=0.5$, $\sigma_\delta=1$\\
Kernel Rasch & as Rasch, $\rho=0.5$, $T=0.1$ & as fitted Rasch, $\rho=0.8$, $T=0.2$ & $m_0=0.048$, $\sigma_\mu=2$, $\sigma_\theta=1$, $\sigma_\delta=1.4$, $\rho=0.8$, $T=0.2$\\
\bottomrule
\end{tabular}
\end{table}

\emph{Fitted on training} means estimated from the responses of the 11 training benchmarks, once, before any simulation starts. For the Beta-Binomial, we choose $(\alpha,\beta)$ to maximise the probability of the training tasks' counts, where a task with $k$ correct labels out of $n$ has probability proportional to $\mathrm{B}(\alpha+k,\beta+n-k)/\mathrm{B}(\alpha,\beta)$, and every benchmark counts equally. This gives a prior mean of $0.522$ worth $\alpha+\beta=3.68$ labels. For Rasch, we first fit the model to $4.34$ million training responses (4,809 subjects, 109,987 items) by maximum a posteriori estimation. We then set each prior parameter from the fitted values: $m_0$ is the mean of the benchmark rates $\mu_d$, and each width is the spread of the corresponding fitted quantity ($\mu_d$, $\theta_i$ or $\delta_j$), averaged per benchmark so that the benchmark with the most subjects does not dominate. We repeat the two steps until the values stop changing. Only these prior parameters transfer to a new benchmark. The fitted $\mu_d$, $\theta_i$ and $\delta_j$ belong to training benchmarks, subjects and items.

\emph{Tuned on validation} is different, because nothing is estimated from a likelihood. We try a grid of settings, run the online simulation on the three validation benchmarks with each, and keep the setting with the lowest Brier ALC. The scale $\gamma$ of the probe prior mean (Section~\ref{sec:probes}) is chosen the same way. The two sources differ in what they learn from and what they optimise. Fitting on training uses many benchmarks and a likelihood and never sees validation or test outcomes. Tuning on validation uses three benchmarks and the score we care about and never sees test outcomes. Both are fixed before the test simulation starts, and the probes are trained on the training benchmarks' items only.

%==============================================================================
\section{Results}\label{sec:results}
%==============================================================================

Every difference is written $\Delta_{a-b}$, the Brier ALC of model $a$ minus that of model $b$, so negative favours $a$; paired 95\% bootstrap intervals are over tasks. J is the Jeffreys Beta-Binomial; BB, R and KR are the Beta-Binomial, Rasch and kernel Rasch models, subscripted fix, tune or fit for the source of their prior, and R\,off is the offline Rasch that keeps its fitted values.

\subsection{Prior source matters, and no source wins on every split}

Table~\ref{tab:prior-source} compares the nine combinations. The validation-tuned Beta-Binomial prior is \emph{worse} than Jeffreys on test ($\D{BB\,tune}{J}=+0.0055$) though it is the best Beta-Binomial on validation. Every comparison between prior sources changes sign between test and validation (Appendix~\ref{app:paired}); fresh draws change no result by more than $0.0007$ (Appendix~\ref{app:robust}). Other metrics tell the same story: log loss and calibration error give the same ordering of the prior sources on test and the same reversal on validation, and rank metrics are uninformative for a constant prediction (Appendix~\ref{app:alternative-metrics}).

\begin{table}[H]
  \caption{\textbf{A prior fitted on training helps on the test benchmarks and hurts on the validation benchmarks; a fixed prior with strong shrinkage does nearly as well on test.} Each cell is $\Delta_{m-\mathrm{J}}$ for the model $m$ in the row, competition-shaped draws. ID: held-out subjects on the 11 training benchmarks, 114 tasks, where the training-fitted Beta-Binomial prior is refit without those subjects. OOD validation: 3 unseen benchmarks, 80 tasks. OOD test: 3 unseen benchmarks, 45 tasks. Bold: best of the nine models in each column; underlined: second best.}\label{tab:prior-source}
  \centering\footnotesize\setlength{\tabcolsep}{0.75pt}
  \begin{tabular}{llccc}
\toprule
Prior source & Model & Test (OOD) & Validation (OOD) & ID \\
\midrule
\multirow{3}{*}{Fixed in advance} & Beta-Binomial (Jeffreys) & reference & reference & reference \\
 & Rasch & $\minus0.0102$ [$\minus0.0118$, $\minus0.0085$] & $+0.0031$ [$\minus0.0005$, $+0.0066$] & $\minus0.0073$ [$\minus0.0091$, $\minus0.0055$] \\
 & Kernel Rasch & $\minus0.0105$ [$\minus0.0121$, $\minus0.0088$] & $+0.0029$ [$\minus0.0008$, $+0.0065$] & $\minus0.0072$ [$\minus0.0087$, $\minus0.0058$] \\
\midrule
\multirow{3}{*}{Tuned on validation} & Beta-Binomial & $+0.0055$ [$+0.0038$, $+0.0072$] & $\minus0.0015$ [$\minus0.0040$, $+0.0010$] & $+0.0041$ [$+0.0027$, $+0.0055$] \\
 & Rasch & $\minus0.0032$ [$\minus0.0035$, $\minus0.0029$] & $+0.0004$ [$\minus0.0006$, $+0.0014$] & $\minus0.0024$ [$\minus0.0028$, $\minus0.0019$] \\
 & Kernel Rasch & $\minus0.0024$ [$\minus0.0027$, $\minus0.0021$] & $\minus0.0003$ [$\minus0.0010$, $+0.0004$] & $+0.0001$ [$\minus0.0004$, $+0.0005$] \\
\midrule
\multirow{3}{*}{Fitted on training} & Beta-Binomial & $\minus0.0138$ [$\minus0.0158$, $\minus0.0116$] & $+0.0079$ [$+0.0024$, $+0.0136$] & $\minus0.0090$ [$\minus0.0115$, $\minus0.0063$] \\
 & Rasch & $\minus0.0119$ [$\minus0.0133$, $\minus0.0106$] & $+0.0052$ [$+0.0008$, $+0.0097$] & $\minus0.0081$ [$\minus0.0101$, $\minus0.0060$] \\
 & Kernel Rasch & $\minus0.0118$ [$\minus0.0131$, $\minus0.0105$] & $+0.0036$ [$\minus0.0007$, $+0.0076$] & $\minus0.0063$ [$\minus0.0076$, $\minus0.0049$] \\
\bottomrule
\end{tabular}

\end{table}

\begin{figure}[t]
  \centering
  \includegraphics{\figdir/figure1-prior-source.pdf}
  \caption{\textbf{The prior's source matters at zero to three labels and has faded by 31.} Brier score (a) and $\Delta_{m-\mathrm{J}}$ (b) per label budget on the three unseen test benchmarks, competition-shaped draws, for the Beta-Binomial with a fixed (Jeffreys, solid), validation-tuned (dashed) and training-fitted (dotted) prior. Jeffreys jumps to $0.75$ or $0.25$ after one answer, so the gap is largest at one label; at 31 labels the three differ by less than $0.001$.}\label{fig:prior-source}
\end{figure}

\subsection{Learned priors transfer to unseen benchmarks; item-level estimates do not}

\begin{figure}[t]
  \centering
  \includegraphics{\figdir/figure2-offline-online.pdf}
  \caption{\textbf{Offline estimates win on familiar benchmarks and have no edge on unfamiliar ones.} Brier score per label budget for the offline Rasch keeping its fitted benchmark rates and item difficulties and updating only ability (solid), against the online Rasch with the same fitted prior (dashed). Left: ID, held-out subjects on the 11 training benchmarks. Right: OOD test benchmarks, where the offline model has no fitted values and falls back to the prior.}\label{fig:offline-online}
\end{figure}

On the ID evaluation set the offline model that keeps its fitted values scores $\D{R\,off}{J}=-0.0852$ [$-0.0955$, $-0.0747$] against $\D{BB\,fit}{J}=-0.0090$ for the fitted Beta-Binomial prior and $\D{R\,fit}{J}=-0.0081$ for the fitted Rasch prior (Figure~\ref{fig:offline-online}); its item difficulties alone, at budget zero, reach a within-task AUC of $0.836$. On the unseen test benchmarks it equals the fitted prior ($\D{R\,off}{J}=-0.0115$ against $\D{R\,fit}{J}=-0.0119$): nothing it knows about particular items or benchmark rates is available. This is the seen-benchmark check the design requires, and it separates two things that are easy to conflate: item-level estimates transfer only to the items they were fitted on, while the prior's spreads carry over to new benchmarks.

\subsection{Learned priors help benchmarks of typical difficulty and hurt much easier or harder ones}

\begin{figure}[t]
  \centering
  \includegraphics{\figdir/figure3-benchmark-breakdown.pdf}
  \caption{\textbf{The training-fitted prior helps on all three test benchmarks and on one validation benchmark, and hurts on the hardest.} $\D{BB\,fit}{J}$ for each unseen benchmark, with 95\% intervals over tasks. The validation benchmark \code{bird\_sql} has task accuracy $0.12$, far below every training benchmark.}\label{fig:benchmarks}
\end{figure}

\begin{figure}[t]
  \centering
  \includegraphics{\figdir/figure4-benchmark-explanation-17.pdf}
  \caption{\textbf{A fitted prior helps when a benchmark's accuracy is near the prior mean.} Each point is one of 17 benchmarks. For each of the 11 training benchmarks, the prior was refit without any of that benchmark's rows and scored on held-out subjects. The six validation and test benchmarks use the original training-only prior. The horizontal axis uses task accuracy on the evaluation items; the vertical axis shows $\D{BB\,fit}{J}$, with 95\% task-bootstrap intervals. This is an explanatory analysis using observed outcomes.}\label{fig:lobo}
\end{figure}

Figure~\ref{fig:benchmarks} shows the test and validation benchmarks separately: all three test benchmarks agree ($\D{BB\,fit}{J}$ from $-0.013$ to $-0.015$), and on validation \code{ai2d\_test} agrees while \code{bird\_sql} goes the other way ($+0.025$). Over the 17 benchmarks (Figure~\ref{fig:lobo}; each training benchmark is held out in turn, Appendix~\ref{app:lobo}) the fitted prior significantly helps on 9, significantly hurts on 1 (\code{bird\_sql}), and the macro mean $\D{BB\,fit}{J}$ is $-0.0050$. What separates the benchmarks is not whether they lie inside the range of training benchmark rates (\code{alignment\_faking} and \code{harmbench} lose while inside it) but how far their task accuracies sit from the prior mean: the Spearman correlation between the gain and that distance is $0.96$ over the 17 benchmarks ($0.93$ over the 11 training benchmarks), against $0.01$ over the 11 for the benchmark's mean accuracy. The prior is a starting guess that a task is answered correctly about $52\%$ of the time, and with few labels the predictor pulls its estimate toward it. If a benchmark's true accuracy is, say, $55\%$, the pull lands near the truth and the prediction improves; if it is $12\%$, as on \code{bird\_sql}, the pull drags the estimate the wrong way and it gets worse. So the prior helps on benchmarks of typical difficulty and hurts on much easier or much harder ones. The three test benchmarks (mean accuracy $0.54$--$0.68$) sit near the prior mean, which is why test looks good. We found this pattern after seeing the results, so it is an explanation and not a predictor.

\par\begin{addnote}{11}A short subsection, \emph{How sensitive is the test gain to which benchmarks are held out?}, moved from the appendix: over 680 triples the mean gain is $-0.0037$ (prior refit on the other 14 benchmarks) or $-0.0051$ (training benchmarks only), the prior hurts in 17.5\%--20\% of splits, and our test triple gives $-0.0138$ when the prior is fitted on the 11 training benchmarks but $-0.0087$ when it is fitted on the other 14. State which fit set is used.\end{addnote}\par
\par\begin{addnote}{12}The first-principles version of the explanation: for a task of true accuracy $p$ the expected Brier excess of the posterior mean is $\frac{(\alpha-p(\alpha+\beta))^2+n\,p(1-p)}{(\alpha+\beta+n)^2}$, so the gain over Jeffreys is an exact function of $p$ and of the budget; compare the predicted gain per benchmark with the observed one instead of citing a post hoc correlation.\end{addnote}\par
Under the fitted priors, the simplest Rasch model is also the best on test and the worst on validation: with benchmark rate only $\Delta_{m-\mathrm{J}}$ is $-0.0172$ on test and $+0.0211$ on validation (Appendix~\ref{app:ladder}). The subject-ability term is what lets a benchmark far from the prior move away from it, at a cost in variance when the benchmark is near it. The kernel adds nothing on test.

\subsection{Item content carries a small rank signal and a negligible gain for the prior}\label{sec:itemfeat}
\label{sec:itemfeat-text}
\begin{table}[H]
\centering
\caption{Item-difficulty transfer: within-benchmark Spearman correlation and pairwise accuracy between predicted easiness and observed pass rate, averaged across benchmarks. ID: items held out within the 11 training benchmarks; OOD validation and OOD test: three unseen benchmarks each. The additional pass-rate and Rasch-target probes use benchmark-grouped selection and were not evaluated on ID (--).}
\label{tab:item-difficulty}
\small\setlength{\tabcolsep}{4pt}
\begin{tabular}{lcccccc}
\toprule
 & \multicolumn{2}{c}{ID} & \multicolumn{2}{c}{OOD validation} & \multicolumn{2}{c}{OOD test}\\
\cmidrule(lr){2-3}\cmidrule(lr){4-5}\cmidrule(lr){6-7}
Probe & Spearman & Pairwise & Spearman & Pairwise & Spearman & Pairwise\\
\midrule
Random score & $-0.003$ & $0.499$ & $-0.024$ & $0.490$ & $0.076$ & $0.529$ \\
Log character count & $0.107$ & $0.539$ & $0.058$ & $0.525$ & $0.061$ & $0.521$ \\
Ridge on embeddings & $0.400$ & $0.649$ & $0.196$ & $0.570$ & $-0.020$ & $0.494$ \\
Pass-rate ridge & -- & -- & $0.165$ & $0.557$ & $-0.035$ & $0.487$ \\
Rasch-target ridge & -- & -- & $0.130$ & $0.545$ & $-0.035$ & $0.487$ \\
Pass-rate MLP & -- & -- & $0.071$ & $0.523$ & $0.109$ & $0.539$ \\
Rasch-target MLP & -- & -- & $0.120$ & $0.540$ & $0.104$ & $0.538$ \\
\bottomrule
\end{tabular}

\end{table}

\begin{qnote}{Q2}Section 5.4: the probe table and the prior-mean table now sit in Results, with the probes defined in Section 4.2. Does the abstract say no more than these two tables show?\end{qnote}
Embeddings show a clear signal on known benchmarks. On the six unseen benchmarks the two MLP probes pass the criterion we fixed in advance, with benchmark-macro Spearman $0.090$ and $0.103$ against $0.026$ for a random baseline and a positive correlation on 5 of 6 benchmarks; the ridge probes do not ($0.065$ and $0.047$). On the three test benchmarks alone, the pass-rate MLP reaches $0.109$, while the Rasch-target MLP reaches $0.059$, below the sampled random control's $0.076$. The aggregate rank criterion therefore does not establish a consistent test-only advantage.

The rank signal yields only a small gain at the validation-selected scale. In Table~\ref{tab:probe-prior}, every $\gamma\ge0.25$ is worse than $\gamma=0$ on both splits. The scale selected on validation ($\gamma=0.05$, Rasch-target MLP) improves Brier ALC by about $0.0004$ on each split in its matched replay; this comparison alone does not establish a reliable improvement. This tests these two probes only; a stronger difficulty predictor, for instance one trained through an item response likelihood on more items, could do better.

\begin{table}[H]
\centering
\caption{Predicted item difficulty as the prior mean of the Rasch item deviation, Equation~\eqref{eq:probe-prior}. Cells are $\Delta_{m-\mathrm{J}}$ of the fitted Rasch with scale $\gamma$, competition-shaped draws; $\gamma=0$ is the fitted Rasch without a prior mean. Bold: lowest in each column.}
\label{tab:probe-prior}
\begin{tabular}{lcccc}
\toprule
 & \multicolumn{2}{c}{Pass-rate MLP} & \multicolumn{2}{c}{Rasch-target MLP} \\
Scale $\gamma$ & OOD validation & OOD test & OOD validation & OOD test \\
\midrule
$0$ & $\mathbf{+0.0052}$ & $\minus0.0119$ & $+0.0052$ & $\minus0.0119$ \\
\midrule
$0.05$ & $+0.0062$ & $\mathbf{\minus0.0119}$ & $\mathbf{+0.0048}$ & $\mathbf{\minus0.0123}$ \\
$0.1$ & $+0.0103$ & $\minus0.0082$ & $+0.0051$ & $\minus0.0120$ \\
$0.25$ & $+0.0332$ & $+0.0150$ & $+0.0099$ & $\minus0.0079$ \\
$0.5$ & $+0.0690$ & $+0.0491$ & $+0.0273$ & $+0.0069$ \\
$1$ & $+0.1065$ & $+0.0826$ & $+0.0704$ & $+0.0433$ \\
$1.5$ & $+0.1326$ & $+0.1050$ & $+0.1040$ & $+0.0724$ \\
\bottomrule
\end{tabular}

\end{table}



\subsection{Priors that adapt to the first labels do not remove the harm}\label{sec:adaptive-results}
Table~\ref{tab:adaptive} compares the adaptive priors of Section~\ref{sec:adaptive} with the fitted Beta-Binomial prior over the 17 benchmarks. They trade gain for safety and none dominates it. The mixtures give the largest average and test gains but leave the worst benchmark as bad as before. Averaging over the mean and strength nearly removes the worst case but halves the gain. The switch rule falls in between. We had fixed three criteria in advance: a worst benchmark of at most $+0.0120$, a mean within $0.0010$ of the fitted prior's $-0.0050$, and a test gain within $0.0020$ of its $-0.0138$. No adaptive prior meets all three.

\begin{table}[H]
\centering
\caption{$\Delta_{m-\mathrm{J}}$ for adaptive priors. Worst benchmark: the largest gain (a loss) over the 17 benchmarks. Mean: over the 17 benchmarks. The training benchmarks are held out one at a time, and the switch threshold is chosen on validation. Mixtures with 2 and 3 components differ by at most $0.0001$ (the table shows 2). Bold: best of the adaptive priors in each column.}
\label{tab:adaptive}
\small
\begin{tabular}{lccc}
\toprule
Model $m$ & Worst benchmark & Mean over 17 benchmarks & OOD test \\
\midrule
Fitted Beta-Binomial (reference) & $+0.0248$ & $\minus0.0050$ & $\minus0.0138$ \\
\midrule
Beta mixture, 2 or 3 components & $+0.0269$ & $\mathbf{\minus0.0054}$ & $\mathbf{\minus0.0163}$ \\
Average over strength & $+0.0122$ & $\minus0.0045$ & $\minus0.0110$ \\
Average over mean and strength & $\mathbf{+0.0034}$ & $\minus0.0037$ & $\minus0.0063$ \\
Switch to Jeffreys, 3 labels & $+0.0136$ & $\minus0.0027$ & $\minus0.0078$ \\
Switch to Jeffreys, 7 labels & $+0.0202$ & $\minus0.0042$ & $\minus0.0114$ \\
\bottomrule
\end{tabular}

\end{table}

Can the first labels warn that the prior will hurt? The distance of the accuracy so far from the prior's mean separates the tasks the prior hurts from the rest, with an area under the ROC curve (AUC) of $0.73$ from seven labels, but only $0.64$ within a benchmark, so most of the signal is between benchmarks. The harm sits at zero to three labels, which carry half of the ALC weight, before any such warning exists.

%==============================================================================
\section{Limitations}\label{sec:limitations}
%==============================================================================
\label{sec:limitations}
The filtered public data covers only three test benchmarks in the online simulation. Multiple draws
and acquisition orders measure variation within this population; they do not
supply independent tests of transfer to new benchmarks. Task bootstrap
intervals likewise do not capture uncertainty from sampling new benchmarks.
Results apply to the eligible response data and random label acquisition,
and need not carry over to other benchmark mixes or selection policies.
The 17-benchmark sensitivity analysis reuses held-out-subject tasks for
the 11 original training benchmarks. The test triples overlap, so their
distribution describes sensitivity within this dataset rather than
independent estimates of transfer.
The ID comparison has item response history that is unavailable under
benchmark cold start. Observed item difficulty also depends on which subjects
answered the item. The leave-one-benchmark-out association uses realized
outcomes to explain when the fitted prior helps; it is not a way to select a
prior before observing a new benchmark.



Further limits of our own. Each competition-shaped draw has one subject per benchmark, so sharing labels across subjects is not tested. These findings concern the models and data studied here, not a general impossibility of item-level transfer. The test gain depends on this split. The test benchmarks are close to the prior mean; over all 680 choices of three held-out benchmarks the fitted prior gains less ($\D{BB\,fit}{J}=-0.0037$ on average, refit on the other 14) and hurts in $17.5\%$ of them (Appendix~\ref{app:split-sensitivity}), as it does on validation. Test scores were inspected repeatedly while building the models, so the test analyses are exploratory: all selection used validation, but test was not read once. The kernel settings were chosen on validation and not on the ID evaluation set. We did not rerun the subject-metadata prior in this design, and we dropped reasoning-trace features, which are not available at test time.

%==============================================================================
\section{Conclusion}
\par\begin{improve}{A-conclusion}Conclusion: still to improve. The adaptive-priors sentence is unsupported until the body reports it; rewrite after the headline and the adaptive-priors result are settled, and keep it ranked and short.\end{improve}\par
%==============================================================================

What transfers from offline fitting to an unseen benchmark is a shrinkage prior, not item-level estimates. Where fitted values apply, offline fitting is worth a very large gain ($\D{R\,off}{J}=-0.0852$ on the ID evaluation set); where they do not, it is worth none. The prior's gain is modest and depends on the benchmark: $-0.0050$ on average over the 17 benchmarks and $-0.0037$ over the 680 choices of three held-out benchmarks, with a loss in $17.5\%$ of them. The prior helps when a benchmark's accuracy lies near the prior mean and hurts when it lies far away, which is why the same fitted prior looked good on test and bad on validation. Whether sharing labels across subjects helps remains untested, since each draw has one subject per benchmark. Two attempts to do better did not help: priors that adapt to the first labels did not remove the harm without losing the gain, and the neural probes we trained on item embeddings, which carry a small rank signal for difficulty, improved prediction by only about $0.0004$ Brier ALC as a prior mean.

%==============================================================================
\section*{Disclosure}
%==============================================================================
\label{sec:disclosure-rerun}
\textbf{Training data.} All fitted quantities (the offline Rasch fit, the Beta-Binomial
empirical-Bayes prior, the item probes) use only the responses of the 11 training benchmarks in
the competition's measurement database after the filtering described in Section~\ref{sec:setting-online};
no other response data is used. \textbf{Pretrained models.} Item text is embedded with the frozen
BGE-M3 model; nothing else is pretrained or fine-tuned.

\textbf{Author contributions and LLM assistance.} The authors formulated the research questions,
directed the experiments, selected the final methods, and interpreted the results. Claude Code
(Anthropic) and OpenAI Codex assisted with brainstorming and refining ideas, implementing code,
preparing analyses and figures, and drafting and revising text. A colleague with training in
statistical learning contributed to discussions of the mathematical formulation. LLM-assisted
development at times became difficult to audit and steer; this was one reason we chose a simpler,
more transparent predictor for the competition submission. All reported numerical results were
generated by executing the repository code and checked by the authors. The authors retained
responsibility for the methodological decisions, verification, and final manuscript.


\clearpage
{\small
\bibliographystyle{plainnat}
\bibliography{references}
}

\appendix

\section{Rasch terms under the fitted priors}\label{app:ladder}
\begin{table}[H]
  \caption{$\Delta_{m-\mathrm{J}}$ for Rasch with terms switched off (prior width $\sigma$ set near zero), under the training-fitted priors. Competition-shaped draws on OOD validation and OOD test. Bold: lowest in each column.}
  \centering\small
  \begin{tabular}{lcc}
\toprule
Model $m$ & OOD validation & OOD test \\
\midrule
Benchmark rate $\mu_d$ only & \begin{tabular}[c]{@{}r@{}}$+0.0211$ \\{}[$+0.0114$, $+0.0311$]\end{tabular} & \begin{tabular}[c]{@{}r@{}}$\minus0.0172$ \\{}[$\minus0.0202$, $\minus0.0142$]\end{tabular} \\
+ subject ability $\theta_i$ & \begin{tabular}[c]{@{}r@{}}$\mathbf{+0.0029}$ \\{}[$\minus0.0003$, $+0.0060$]\end{tabular} & \begin{tabular}[c]{@{}r@{}}$\minus0.0089$ \\{}[$\minus0.0100$, $\minus0.0080$]\end{tabular} \\
+ item deviation $\delta_j$ (Rasch) & \begin{tabular}[c]{@{}r@{}}$+0.0052$ \\{}[$+0.0007$, $+0.0098$]\end{tabular} & \begin{tabular}[c]{@{}r@{}}$\minus0.0119$ \\{}[$\minus0.0133$, $\minus0.0106$]\end{tabular} \\
+ kernel & \begin{tabular}[c]{@{}r@{}}$+0.0036$ \\{}[$\minus0.0006$, $+0.0078$]\end{tabular} & \begin{tabular}[c]{@{}r@{}}$\minus0.0118$ \\{}[$\minus0.0131$, $\minus0.0105$]\end{tabular} \\
\midrule
$\mu_d$ and $\delta_j$, no $\theta_i$ & \begin{tabular}[c]{@{}r@{}}$+0.0282$ \\{}[$+0.0164$, $+0.0401$]\end{tabular} & \begin{tabular}[c]{@{}r@{}}$\mathbf{\minus0.0174}$ \\{}[$\minus0.0208$, $\minus0.0139$]\end{tabular} \\
$\theta_i$ and $\delta_j$, no $\mu_d$ & \begin{tabular}[c]{@{}r@{}}$+0.0078$ \\{}[$+0.0022$, $+0.0135$]\end{tabular} & \begin{tabular}[c]{@{}r@{}}$\minus0.0140$ \\{}[$\minus0.0157$, $\minus0.0122$]\end{tabular} \\
\bottomrule
\end{tabular}

\end{table}

\section{Paired comparisons between prior sources}\label{app:paired}
\begin{table}[H]
  \caption{Paired $\Delta_{a-b}$ between prior sources, with 95\% intervals over tasks. Each comparison changes sign between OOD test and OOD validation. A dash: not run on ID.}
  \centering\small\setlength{\tabcolsep}{4.5pt}
  \begin{tabular}{lccc}
\toprule
Comparison (a $-$ b) & Test & Validation & ID \\
\midrule
Rasch fitted $-$ Rasch fixed & $\minus0.0017$ [$\minus0.0028$, $\minus0.0005$] & $+0.0021$ [$+0.0006$, $+0.0036$] & -- \\
Kernel fitted $-$ Rasch fitted & $+0.0001$ [$\minus0.0002$, $+0.0004$] & $\minus0.0016$ [$\minus0.0021$, $\minus0.0012$] & -- \\
Rasch fitted $-$ Beta-Binomial fitted & $+0.0019$ [$+0.0007$, $+0.0031$] & $\minus0.0027$ [$\minus0.0041$, $\minus0.0012$] & $+0.0009$ [$\minus0.0001$, $+0.0018$] \\
\bottomrule
\end{tabular}

\end{table}

\section{Draw-seed robustness}\label{app:robust}
\begin{table}[H]
  \caption{$\Delta_{m-\mathrm{J}}$ on three fresh sets of 150 competition-shaped draws per role (seeds 20261201--03); results on the frozen draws are in Table~\ref{tab:prior-source}.}
  \centering\small
  \begin{tabular}{lcc}
\toprule
Model & Test & Validation \\
\midrule
Beta-Binomial fitted & $\minus0.0137$ / $\minus0.0144$ / $\minus0.0139$ & $+0.0078$ / $+0.0083$ / $+0.0089$ \\
Rasch fitted & $\minus0.0120$ / $\minus0.0124$ / $\minus0.0123$ & $+0.0051$ / $+0.0054$ / $+0.0060$ \\
Beta-Binomial tuned & $+0.0060$ / $+0.0059$ / $+0.0063$ & $\minus0.0015$ / $\minus0.0016$ / $\minus0.0018$ \\
\bottomrule
\end{tabular}

\end{table}

\section{The fitted prior on all 17 benchmarks}\label{app:lobo}
\begin{table}[H]
  \caption{$\D{BB\,fit}{J}$ for each of the 17 benchmarks. Training benchmarks: the prior is refit without the held-out benchmark and scored on held-out-subject tasks. OOD validation and OOD test benchmarks: the original training-only prior. The last column is the mean distance of a benchmark's task accuracies from the prior mean $0.52$.}
  \centering\small
  \begin{tabular}{lccc}
\toprule
Held-out benchmark & Tasks & $\Delta_{\mathrm{BB\,fit}-\mathrm{J}}$ [95\% CI] & Mean distance from $0.52$ \\
\midrule
\code{hallusionbench} & 20 & $\minus0.0151$ [$\minus0.0165$, $\minus0.0136$] & 0.11 \\
\code{bridging\_gap} & 20 & $\minus0.0142$ [$\minus0.0167$, $\minus0.0114$] & 0.14 \\
\code{helm\_afr} & 3 & $\minus0.0139$ [$\minus0.0168$, $\minus0.0108$] & 0.07 \\
\code{afrimedqa} & 6 & $\minus0.0121$ [$\minus0.0189$, $\minus0.0046$] & 0.15 \\
\code{afrieval} & 20 & $\minus0.0111$ [$\minus0.0138$, $\minus0.0085$] & 0.18 \\
\code{cruxeval} & 5 & $\minus0.0055$ [$\minus0.0155$, $+0.0056$] & 0.17 \\
\code{brace} & 6 & $\minus0.0042$ [$\minus0.0121$, $+0.0037$] & 0.22 \\
\code{agi\_elo} & 4 & $+0.0004$ [$\minus0.0069$, $+0.0076$] & 0.26 \\
\code{appworld} & 4 & $+0.0026$ [$\minus0.0087$, $+0.0120$] & 0.24 \\
\code{alignment\_faking} & 20 & $+0.0043$ [$\minus0.0092$, $+0.0188$] & 0.24 \\
\code{harmbench} & 6 & $+0.0070$ [$\minus0.0023$, $+0.0161$] & 0.31 \\
\midrule
Mean over the 11 benchmarks & 114 & $\minus0.0056$ & 0.19 \\
\bottomrule
\end{tabular}

\end{table}

\par\begin{addnote}{7}New appendix tables: the adaptive-prior results per benchmark; the pilot-diagnostic AUC and the switch-rule sweep; and the per-split distribution for the Rasch prior (40 triples).\end{addnote}\par
\section{Sensitivity to the benchmark split}\label{app:split-sensitivity}
We held out each of the 680 triples, refitted the benchmark-weighted empirical Beta prior on the other 14 benchmarks, and rescored the frozen trajectories. The mean macro $\Delta_{\mathrm{BB\,fit}-\mathrm{J}}$ was $\minus0.0037$, with 5th--95th percentiles [$\minus0.0095$, $+0.0099$]. The fitted prior hurt on 17.5\% of splits. The original test benchmark names gave $\minus0.0087$ (percentile 13) under refitting. This refit includes the former validation benchmarks among the 14 eligible fit benchmarks for the original test triple, so its result differs from the original 11-benchmark fit. The 11 original training benchmarks are scored on held-out-subject tasks from the existing leave-one-out data. These overlapping triples describe split sensitivity within this dataset. The main text uses the prior fitted on the 11 training benchmarks only. Under that prior the mean macro $\Delta_{\mathrm{BB\,fit}-\mathrm{J}}$ over the same triples is $\minus0.0050$, the prior hurts on 17.8\% of them, and the original test benchmarks give $\minus0.0138$.

\begin{figure}[H]
  \centering
  \includegraphics{\figdir/figure5-split-sensitivity.pdf}
  \caption{$\D{BB\,fit}{J}$ over all 680 triples. Each empirical prior is fitted on the 14 benchmarks outside its triple. The vertical line marks the original test benchmark names under this refitting protocol.}
\end{figure}

\section{Alternative evaluation metrics}\label{app:alternative-metrics}
\begin{table}[H]
  \centering\small
  \caption{Competition-shaped online simulation; models are labelled as in Section~\ref{sec:results}. Brier, log loss, and ten-bin task ECE are averaged using the official budget weights; lower is better. AUC uses within-task item rankings at 31 labels, where higher is better. Dashes indicate no informative ranking; a constant prediction ties all items. Bold: best in each split (lowest loss, highest AUC). Per-budget AUC and alternative weights are reported in the accompanying CSVs.}
  \begin{tabular}{lrrr}
\toprule
Model & Brier ALC & Log loss ALC & ECE ALC \\
\midrule
Beta-Binomial (Jeffreys) & 0.2499 & 0.7045 & 0.1426 \\
Beta-Binomial (tuned) & 0.2555 & 0.7162 & 0.1522 \\
Beta-Binomial (fitted) & 0.2362 & 0.6656 & 0.1113 \\
Rasch (fixed) & 0.2397 & 0.6739 & 0.1203 \\
Rasch (tuned) & 0.2467 & 0.6940 & 0.1358 \\
Rasch (fitted) & 0.2380 & 0.6712 & 0.1164 \\
Kernel Rasch (fixed) & 0.2394 & 0.6733 & 0.1218 \\
Kernel Rasch (tuned) & 0.2475 & 0.6972 & 0.1428 \\
Kernel Rasch (fitted) & 0.2382 & 0.6717 & 0.1230 \\
\bottomrule
\end{tabular}

\end{table}

\section{Solver convergence}\label{app:convergence}
\paragraph{Online solver.} At every budget the online Rasch fit is a Newton solve with a limit of 30 steps. In 6,300 audited fits, none reached the limit, none needed more than 7 steps, and the last step was below $10^{-6}$ in every one.

\paragraph{Offline fit.} Our first offline fit stopped at its iteration limit of 400 in three of its four rounds. We refit it with a proper stopping rule. The hyperparameters moved by at most $1.2\%$, the benchmark rates by at most $0.011$ on the logit scale, and no Brier ALC changed by more than $0.0001$. The reported results use the first fit, and the refit leaves them unchanged.

\section{Probe details}\label{app:probes}
\paragraph{Training.} Inputs are standardised using the fitting items only. The neural network has one hidden layer of 32 units and a limit of 100 epochs. The Rasch targets come from the offline Rasch fit.

\paragraph{From predictions to $z_j$.} For the Rasch-target probe, the difficulty $z_j$ is the negative of the predicted easiness. For the pass-rate probe, the centred prediction $f(x_j)$ is converted to a logit-scale difficulty,
\begin{equation*}
z_j=\operatorname{logit}(\bar p)-\operatorname{logit}\!\big(\operatorname{clip}(\bar p+f(x_j),\,10^{-4},\,1-10^{-4})\big),
\end{equation*}
where $\bar p$ is the benchmark-weighted mean pass rate of the sampled training items. Positive difficulty means harder.


\clearpage
\section*{What I would change and why}
\small
\begin{enumerate}\setlength{\itemsep}{2pt}
\item \textbf{Headline.} $0.014$ is the best case: our test split is one of the most favourable we could have drawn. The average over splits is about $0.004$--$0.005$, and the prior hurts in roughly one split in five.
\item \textbf{17 benchmarks.} Figure~4 already shows 17 benchmarks (Spearman $0.96$, helps on 9, hurts on 1); the text, the abstract and the appendix table still say 11 (Spearman $0.93$, 5 of 11).
\item \textbf{Item features.} A pass-rate MLP and a Rasch-target MLP pass the rank criterion set in advance (pooled $0.090$ and $0.112$ against $0.026$ for a random control); ridge probes do not. On the three test benchmarks alone the margin over the control is small ($0.109$ against $0.076$). The claim of no signal must be softened; the Brier test of the prior-mean option decides whether it matters.
\item \textbf{Label sharing.} ``Adds nothing to a good prior'' holds on test and with one subject per benchmark. On validation the fitted Rasch beats the fitted Beta-Binomial, and multi-subject draws are barely tested.
\item \textbf{Negative results the reviewer asked for.} Adaptive priors and the pilot switch rule do not meet the criteria fixed in advance; the first labels warn only weakly.
\item \textbf{Specificity.} Say which widths and kernel settings the validation-tuned Rasch priors took.
\item \textbf{Convergence and extra appendix tables} (answers a highlighted reviewer point).
\item \textbf{Wording of the best result.} Add that it is the favourable split and that the validation split reverses.
\item \textbf{Other metrics.} Log loss and calibration agree with Brier (reviewer question 3).
\item \textbf{Correct the 11-benchmark numbers} in the running text to the 17-benchmark view.
\item \textbf{Sensitivity to the split in the main text,} with the fit set stated: $-0.0138$ (11 training benchmarks) against $-0.0087$ (the other 14).
\item \textbf{An exact risk formula} so the explanation is a derivation, not a correlation found after the fact.
\item \textbf{Limitations.} Quantify the favourable-split statement with the 680 splits.
\item \textbf{Conclusion.} Rewrite around the average, the spread across splits and the limits of what a label-based diagnostic can do.
\item \textbf{Terminology.} Done: the table caption now says ``online simulation'' (see S5).
\item \textbf{Appendix table} to 17 benchmarks.
\end{enumerate}
\section*{Edits already made to the paper (old wording red, new wording green)}
\begin{enumerate}\setlength{\itemsep}{2pt}
\item[\textbf{S1}] \textbf{Beta-Binomial notation.} Conventional symbols (checked against standard references): success probability $p$, prior $\mathrm{Beta}(\alpha,\beta)$, $n$ trials with $k$ successes, posterior $\mathrm{Beta}(\alpha+k,\beta+n-k)$ and posterior mean $(\alpha+k)/(\alpha+\beta+n)$, with $\alpha+\beta$ read as the number of pseudo-labels the prior is worth. The old text used $\pi$, a mean $m$ and a strength $s$ (which clashed with the subject index $s$) and $h$ for the correct labels. The prior settings are now stated as $(\alpha,\beta)$: Jeffreys $(\tfrac12,\tfrac12)$, validation-tuned $(0.4,0.6)$, training-fitted $(1.92,1.76)$.
\item[\textbf{S2}] \textbf{Rasch equation in IRT convention.} $p_{ij}=\sigma(\theta_i-b_j)$, replacing $\operatorname{logit}p=\mu_b+\theta_s-b_i$. The item difficulty is $b_j$ (see S4 for its split); the IRT symbol $\beta$ is left to the Beta prior.
\item[\textbf{S3}] \textbf{Response matrix.} Section 2 now defines $Y\in\{0,1\}^{N\times M}$ with rows $i$ (subjects) and columns $j$ (items), benchmarks as blocks of columns, a task as one row restricted to one block, and the splits as disjoint column blocks; the ID control holds out rows.
\item[\textbf{S4}] \textbf{Symbols.} Item difficulty $b_j=\delta_j-\mu_{d(j)}$: $\delta_j$ is the deviation from the benchmark's average, $\mu_d$ the benchmark rate, $d(j)$ the dataset (benchmark) of item $j$, $D=17$ benchmarks. Prior widths are standard deviations $\sigma_\mu,\sigma_\theta,\sigma_\delta$ (formerly $\tau$); the kernel uses an indicator $\mathbf 1_{j=j'}$.
\item[\textbf{S6}] \textbf{Explicit differences.} Every difference is $\Delta_{a-b}$, the Brier ALC of model $a$ minus model $b$, with J (Jeffreys), BB, R and KR (Beta-Binomial, Rasch, kernel Rasch) and fix, tune, fit or off as the subscript, in the Results text, table headers, the paired and robustness tables and the captions. \textbf{Not yet done:} the axis labels inside Person-B's figure files still read ``difference to Jeffreys'' or ``Fitted minus Jeffreys''.
\item[\textbf{S7}] \textbf{Related work} is a numbered section after the Introduction instead of a bold run-in paragraph.
\item[\textbf{S8}] \textbf{Title block.} Authors, affiliation, address and emails added (were \emph{Anonymous Author(s)} and the repository link). The page footer keeps the NeurIPS 2026 conference line.
\item[\textbf{S9}] \textbf{Methodology.} Section 4, formerly \emph{Models and prior sources}, is now \emph{Methodology} with two subsections: \emph{Models} (Beta-Binomial, Rasch, kernel Rasch as its own paragraph, the offline model) and \emph{Priors} (a sentence listing each model's prior parameters, then the three sources: fixed, tuned, fitted).
\item[\textbf{S10}] \textbf{Appendix.} Same symbols as the main text: rows of the Rasch-terms table use $\mu_d,\theta_i,\delta_j$ and the width $\sigma$; table rows and the paired, robustness and metrics tables use the model labels J, BB, R, KR (fix, tune, fit, off) with $\Delta_{a-b}$; ``panel'' removed.
\item[\textbf{S11}] \textbf{Section 3 structure.} The protocol text is split into bold run-in paragraphs: Data, Splits, Draws, Label acquisition, Scoring, In-distribution control. The wording is unchanged apart from S3.
\item[\textbf{S12}] \textbf{Kernel Rasch.} The two item priors are shown side by side as equations: Rasch $\boldsymbol\delta\sim\mathcal N(\mathbf 0,\sigma_\delta^2 I)$ and kernel Rasch $\boldsymbol\delta\sim\mathcal N(\mathbf 0,K)$ with $K=\sigma_\delta^2[(1-\rho)I+\rho S]$, so the only difference is visible; they coincide at $\rho=0$.
\item[\textbf{S13}] \textbf{Abstract.} ``language model'' is now ``LLM''; the rest of the abstract is unchanged, so its claims still await [1], [3] and [4].
\item[\textbf{S14}] \textbf{Abstract, authors.} Added the 17-benchmark average and the 680-split average with the share of harmful splits (the favourable-split point of [1]); replaced the item-features sentence with the probe result and the R5 null; added the equal-contribution mark (asterisk on both names, footnote ``Work done equally.'').
\item[\textbf{S15}] \textbf{Locked in.} Items [2], [3], [10], [13] and [14] are now in the paper and no longer shown as red or green: the 17-benchmark view (Spearman $0.96$, helps on 9 of 17, hurts on 1, macro $-0.0050$) in the abstract, the contributions and Section 5.3; the two neural probes and the R5 null in Section 6; the 680-split sentence in the limitations; the rewritten conclusion.
\item[\textbf{Q1, Q2}] \textbf{Open questions} on the 17-benchmark consistency and on Section 6 are shown as green question boxes in the text.
\item[\textbf{S16}] \textbf{Embeddings and kernel settings.} A new \emph{Item embeddings} paragraph (frozen BGE-M3, 1024 dimensions, 8,192-token context, unit length) is in Section 3 and Kernel Rasch refers to it; Section 3 is four bold paragraphs (Data, Item embeddings, Splits and draws, Label acquisition and scoring); the kernel paragraph states how unlabelled items are predicted; the settings moved to \emph{Priors}: tuned values for Rasch and kernel Rasch, and that the fitted kernel Rasch keeps $\rho=0.8$, $T=0.2$ from validation.
\item[\textbf{S17}] \textbf{Honest claims.} Abstract and conclusion lead with the average over 17 benchmarks and over the 680 splits and the share of harmful splits, then give the test and validation gains as the two examples. The label-sharing claim is replaced by a statement that it is untested (one subject per benchmark in every draw), also added to the limitations. The item-feature claim is limited to the two probes we trained (Section 6 notes a stronger predictor could do better).
\item[\textbf{S18}] \textbf{Item-difficulty section dissolved.} The separate Section 6 is gone. \emph{Probes} is a new subsection of Methodology (4.2: the probes, how they are scored, and Equation 5, the probe prior mean); the two tables and the discussion are Results 5.4, \emph{Do item features help?}; the scale $\gamma$ and the training-only probes are noted in \emph{Priors}; a contribution bullet was added.
\item[\textbf{S19}] \textbf{Research questions and headings.} The Results subsections are statements of what each shows (no separate list of research questions: the Introduction keeps its one governing question and the contribution bullets). Table headers read ID, OOD validation, OOD test (Table 1 reordered to match; Table 2 caption gives the benchmark counts).
\item[\textbf{F1, F3, F4, F5}] \textbf{Figure labels: done by Person-B (commit 2bea768).} The axis labels and legends inside Figures 1, 3, 4 and 5 still use the old wording (``difference to Jeffreys'', ``Fitted minus Jeffreys'', ``Test'', ``Validation''). They are drawn by Person-B's plot scripts, so Person-B must regenerate them with $\Delta_{a-b}$ labels and OOD legends; shown as red boxes at each figure.
\item[\textbf{S20}] \textbf{Tables.} All tables follow NeurIPS style (booktabs, captions above, no vertical rules). Table 1 and Tables 3, 4 and 8 bold the best value per column with a note in the caption; Table 1 shows the reference as ``0 (reference)''; Table 2 has grouped columns (Spearman, Pairwise under ID, OOD validation, OOD test) with the random score as a baseline row; Table 4 separates the added terms from the ablations; Table 7 has a mean row; Table 8 has a rule between splits and the split label on the first row of each block.
\item[\textbf{S21}] \textbf{Person-B's update incorporated.} Original-fit probe numbers (Rasch-target test Spearman $0.059$, below the random control's $0.076$), Table 2 rebuilt from the confirmed CSVs, and the R5 wording softened: the validation-selected scale improves Brier ALC by only about $0.0004$ on each split. The abstract, the contribution bullet, the Section 5.4 heading and the conclusion now say this instead of ``no scale helps''.
\item[\textbf{S5}] \textbf{Wording.} Section 2 is titled \emph{Data and evaluation protocol} (was \emph{Setting and online simulation}); ``panel'' became ``ID control set'' and ``response data''; ``replay'' became ``online simulation''.
\end{enumerate}
\end{document}
